% Paper D position draft. Not a results paper; not submitted.
\documentclass[11pt]{article}
\pdfoutput=1
\usepackage[T1]{fontenc}
\usepackage[utf8]{inputenc}
\usepackage{lmodern}
\usepackage[margin=1in]{geometry}
\usepackage{amsmath,amssymb}
\usepackage{booktabs,longtable,array}
\usepackage{listings}
\usepackage{xcolor}
\usepackage[hyphens]{url}
\usepackage[colorlinks=true,linkcolor=black,citecolor=black,urlcolor=blue!50!black]{hyperref}
\usepackage{microtype}
\setlength{\parskip}{0.5em}
\linespread{1.08}
\lstset{basicstyle=\ttfamily\footnotesize,breaklines=true,frame=single,
  columns=fullflexible,keepspaces=true,showstringspaces=false}

\title{Module Boundaries and Gradual Memory Safety on a C99-Compatible\\
Compiler Architecture}
\author{Yunzuo Chen\\
PartnerNet Software Pty Ltd, Australia\\
\texttt{wanjochan@partnernetsoftware.com}}
\date{Position draft, October 2026}

\begin{document}
\maketitle

\begin{abstract}
This position paper describes the \texttt{.cx} architecture: a POSIX C99-compatible
source surface with explicit module declarations, checked public interfaces, and a
build graph that can be audited without relying on implicit \texttt{\#include}
coupling. The architecture is intended to serve three execution backends from one
module contract: the native \texttt{unisacc} compiler and the host drivers
\texttt{gcc\_cx}/\texttt{clang\_cx}. We separate module boundaries from memory
safety, and separate deterministic checks from optional finite-network inference.
The paper is an intent and design record, not a claim that Rust-style compile-time
ownership, complete memory safety, or release qualification has been achieved.
\end{abstract}

\section{Position and non-goals}

The target is not a new language. The source remains POSIX C99; \texttt{.cx} adds
module, import, public-set, and package declarations through a pragma/config
surface. The host compilers ignore unknown pragmas; expansion belongs to the
driver. The native compiler and the host drivers must agree on module visibility,
refusal diagnostics, and the build graph.

Non-goals of this document are: (i) claiming that runtime handle checks are
compile-time ownership; (ii) claiming a soundness proof that does not exist;
(iii) replacing the Paper~A stream/network pipeline; (iv) treating a green local
test as release qualification.

\section{Two layers}

\begin{longtable}{@{}p{0.22\textwidth}p{0.35\textwidth}p{0.35\textwidth}@{}}
\toprule
Layer & Claim & Evidence direction \\
\midrule
\textbf{A: \texttt{.cx} architecture} &
Module/use graph, exported and private names, explicit resources, build closure,
and staged ownership/borrowing contracts. &
\texttt{cx/}, \texttt{tests/cx*.cx}, \texttt{plans/cx-modules.md}; native and host
drivers must share one contract. \\
\textbf{B: classic C99 optional node} &
Without editing source lines: model-assisted annotation, deterministic checker,
and graded compile/tape fallback. &
Separate v0.1.3 study; not the main \texttt{.cx} claim. \\
\bottomrule
\end{longtable}

Module boundaries are a precondition for a safety contract, not the contract
itself. A clear boundary does not prove the absence of use-after-free or
out-of-bounds access.

\section{Division of labour}

\begin{longtable}{@{}p{0.27\textwidth}p{0.31\textwidth}p{0.32\textwidth}@{}}
\toprule
Work & Owner & Reason \\
\midrule
Module graph and exports & deterministic compiler/driver & Parsing, public sets,
package manifests, cache invalidation; dependencies are not guessed by weights. \\
Optional shape/ownership inference & finite key$\to$class network & Same shape
as Paper~A tables; a wrong inference may only request a runtime check or a named
refusal, never turn an unproved case into ``proved''. \\
Safety-contract validation & deterministic checker & The source of the guarantee;
a table/network cannot replace soundness of the checking rules. \\
Unproved fallback & runtime check or refusal & Unproved input is either checked or
explicitly rejected; strict mode emits no false receipt. \\
\bottomrule
\end{longtable}

The goal is zero false negatives under an explicitly declared language subset and
threat model. That proof does not yet exist. False positives and runtime overhead
are measurable, versioned metrics, not safety claims. Runtime handle isolation
(e.g.\ the \texttt{owned} prototype) and compile-time Rust-style ownership are
separate qualifications.

\section{Staged deliverables}

\begin{enumerate}
\item \textbf{Verifiable module surface.} No new module \texttt{\#include}
  ratchet; package manifest and CXIF; explicit build closure; native and host
  drivers reject the same malformed module graphs.
\item \textbf{Minimal safety contract.} Use-after-release, out-of-bounds, handle
  generation/drop; positive and negative fixtures; a ratchet that only decreases.
  Start with runtime constraints, then a compile-time checking slice.
\item \textbf{Optional inference node.} Annotate only declared contracts; compare
  the provable fraction and runtime overhead against unchanged C99, Lua, SQLite,
  and miniz vectors if the classic-C layer is studied.
\item \textbf{Consistent grading.} The same function has the same observable
  semantics under ``proved'', ``runtime-checked'', and ``refused''; bypassing a
  check cannot produce a valid receipt.
\end{enumerate}

\section{Theory questions}

\begin{itemize}
\item Module boundaries: export set, cross-module calls, and
  \texttt{static}/manifest consistency.
\item Soundness of annotations and borrow rules: checker acceptance implies no
  error in the declared threat model, with an explicit assumption list.
\item Inference domain: is the key space finite? How are tables split when it is
  too large, following the Paper~A construction discipline?
\item Graded compilation: correctness and non-forgeable receipts for each mode.
\end{itemize}

\section{Experiments and reproducibility}

The CX evidence should be generated by gates and referenced, not hand-copied:
ratchet counts, package/CXIF fixtures, ownership red-case matrices, and
native-vs-host parity. A single green fixture does not imply complete memory
safety. Every backend experiment must record compiler identity, input bytes, cold/hot
conditions, and the failure/refusal domain.

\section{Related work and risks}

Rust rewrites the source and uses compile-time borrow checking. Checked~C and
Cyclone rely on manual annotations. Fil-C, CHERI, and tape-safety modes are
fallback mechanisms; a proved path should return to native speed. The primary
risks are (i) presenting runtime prototypes as compile-time qualification,
(ii) presenting module progress as memory-safety progress, (iii) an inference
domain that is not actually finite, and (iv) alias/VLA assumptions expanding
without bound.

\section{Index}

\begin{itemize}
\item \texttt{cx/README.md}: structured CX entry.
\item \texttt{docs/cx/prd.md}: current product reading.
\item \texttt{plans/cx-modules.md}: module and safety ratchet stages.
\item \texttt{archive/research-2026-11d/plans/v0.0.44.md}: C44-CX-MODULE and C44-CX-MEMORY registration.
\item \texttt{docs/cx-spec-draft.md}: language draft.
\item \texttt{research/paper-c-intent.md}: pipeline method.
\item \texttt{research/unisacc-paper.md}: Paper~A, stream/network pipeline.
\end{itemize}

\section*{Status}

This file is a position draft generated from
\texttt{research/paper-d-intent.md}. It contains no measured safety result, no
soundness proof, and no release claim. The Markdown intent remains the source of
truth until the paper is independently reviewed.
\end{document}
